\documentclass[11pt]{article}

\usepackage[T1]{fontenc}
\usepackage{amsmath,amssymb,amsthm}
\usepackage{microtype}
\usepackage[hidelinks]{hyperref}
\usepackage[margin=1.15in]{geometry}

\newtheorem{theorem}{Theorem}
\newtheorem{lemma}[theorem]{Lemma}
\newcommand{\RR}{\mathbb{R}}
\newcommand{\La}{\mathcal{L}_a}

\title{Multiplication Closure with a Linear Factor\\
for Normalized Pochhammer Transforms}
\author{Xiaojian Zeng\\
\small PBC School of Finance, Tsinghua University, 100190\\
\small \texttt{xiaojianzeng@foxmail.com}}
\date{}

\begin{document}
\maketitle

\begin{abstract}
Fix $a>0$, and let $\mathcal L_a$ send the normalized falling factorial
$x^{\underline{k}}/(a)_k$ to $x^k$.  We prove that if
$\mathcal L_a(P_1)$ has degree one and all its zeros in $[-1,0]$, while
$\mathcal L_a(P_2)$ has arbitrary finite degree and all its zeros in the
same interval, then every zero of $\mathcal L_a(P_1P_2)$ also belongs to
$[-1,0]$.  In particular, this resolves the degree-$(1,2)$ case of the
all-$a$ multiplication-closure question.  The proof identifies multiplication
by the linear factor with a first-order differential operator and locates its
zeros by a strictly increasing partial-fraction quotient.  Repeated and
endpoint zeros follow by a degree-preserving compact limit.
\end{abstract}

\section{Context and statement}

For $k\geq 0$, write
\[
 (a)_0=1,\qquad (a)_k=a(a+1)\cdots(a+k-1),\qquad
 e_k(x)=\frac{x^{\underline{k}}}{(a)_k},
\]
where $x^{\underline{k}}=x(x-1)\cdots(x-k+1)$.  The polynomials
$e_k$ form a basis of $\RR[x]$.  Define the degree-preserving linear
isomorphism
\[
 \La:\RR[x]\longrightarrow\RR[x],\qquad \La(e_k)=x^k.
\]
This is the normalized Pochhammer transform introduced in the present form
by Vishnyakova~\cite{Vishnyakova2026}.  It belongs to the broader study of
linear operators preserving real-rootedness; general structural results for
such preservers are given, for example, by Borcea and
Brändén~\cite{BorceaBranden2009}.

Vishnyakova~\cite[Definition~1.4 and Open Problem~1.9]{Vishnyakova2026}
asks whether the class of polynomials whose transforms have all zeros in
$[-1,0]$ is closed under multiplication for every $a>0$.  The same source
records Wagner's theorem that this closure holds when $a=1$; Wagner's
original result is in~\cite{Wagner1992}.  The theorem below proves the full
all-$a$ assertion whenever one factor has transformed degree one.  It does
not address the case in which both transformed factors have degree at least
two.

Exact rational certificates are another useful way to settle restricted
parameter families in operator-polynomial problems; a recent example is
Zeng's treatment of quadratic families in a spectral-set setting
\cite{Zeng2026Lens}.  Here, however, the universal parameter range is handled
directly by monotonicity, and computation is not part of the proof.

\begin{theorem}[Linear-factor multiplication closure]\label{thm:main}
Let $a>0$, let $n\geq0$, and let $P_1,P_2\in\RR[x]$ be nonzero polynomials
with $\operatorname{deg} P_1=1$ and $\operatorname{deg} P_2=n$.  If every zero of $\La(P_1)$ and
$\La(P_2)$ lies in $[-1,0]$, then every zero of $\La(P_1P_2)$ lies in
$[-1,0]$.
\end{theorem}

Taking $n=2$ gives the first nontrivial low-degree instance of the
multiplication-closure question for every $a>0$.

\section{The differential operator}

The key reduction uses only the elementary falling-factorial identity
\[
 x\,x^{\underline{k}}=x^{\underline{k+1}}+k x^{\underline{k}}.
\]
Since $e_1=x/a$, it follows that
\begin{equation}\label{eq:basis-product}
 e_1e_k=\frac{a+k}{a}e_{k+1}+\frac{k}{a}e_k.
\end{equation}

\begin{lemma}\label{lem:operator}
Let $f=\La(P)$.  For every $r\in\RR$,
\begin{equation}\label{eq:operator}
 \La\bigl((e_1+r)P\bigr)
   =(x+r)f(x)+\frac{x(x+1)}{a}f'(x).
\end{equation}
\end{lemma}

\begin{proof}
Write $f(x)=\sum_{k=0}^n b_kx^k$, so that
$P=\sum_{k=0}^n b_ke_k$.  Applying $\La$ to
\eqref{eq:basis-product} coefficientwise gives
\[
 \La(e_1P)=\sum_{k=0}^n b_k
 \left(\frac{a+k}{a}x^{k+1}+\frac{k}{a}x^k\right)
 =xf+\frac{x(x+1)}a f'.
\]
Adding $r\La(P)$ proves the identity.
\end{proof}

Thus it remains to study, for $0\leq r\leq1$, the operator
\begin{equation}\label{eq:Tr}
 T_rf=(x+r)f+\frac{x(x+1)}a f'.
\end{equation}

\section{Root location}

\begin{lemma}\label{lem:location}
Fix $a>0$ and $0\leq r\leq1$.  If every zero of a nonzero real polynomial
$f$ lies in $[-1,0]$, then every zero of $T_rf$ lies in $[-1,0]$.
Moreover, if $f$ has degree $n$, then $T_rf$ has degree $n+1$.
\end{lemma}

\begin{proof}
First suppose that $n\geq1$, that $f$ is monic with distinct zeros in
$(-1,0)$, and that $0<r<1$.  Write
\[
 f(x)=\prod_{i=1}^n(x+s_i),\qquad 0<s_1<\cdots<s_n<1.
\]
Away from the zeros of $f$, logarithmic differentiation and elementary
division give
\begin{align}
 H(x):=\frac{T_rf(x)}{f(x)}
 &=x+r+\frac{x(x+1)}a\sum_{i=1}^n\frac1{x+s_i}\notag\\
 &=\left(1+\frac na\right)x+r+\frac{n-\sum_i s_i}{a}
   -\sum_{i=1}^n\frac{s_i(1-s_i)}{a(x+s_i)}.\label{eq:H}
\end{align}
Consequently, on every interval separated by the poles $-s_i$,
\begin{equation}\label{eq:Hprime}
 H'(x)=1+\frac na+
 \sum_{i=1}^n\frac{s_i(1-s_i)}{a(x+s_i)^2}>0.
\end{equation}
Each pole has strictly negative residue.  Hence $H$ tends to $+\infty$
from the left of each pole and to $-\infty$ from its right.  Together with
the signs at infinity and strict monotonicity, this puts exactly one zero of
$H$ in each of the $n+1$ pole-separated intervals.  No polynomial zero is
lost at a pole, because
\[
 T_rf(-s_i)=-\frac{s_i(1-s_i)}a f'(-s_i)\neq0.
\]

The two exterior zeros also lie in the target interval.  Indeed, directly
from the first line of \eqref{eq:H},
\begin{equation}\label{eq:endpoint-values}
 H(-1)=r-1<0,\qquad H(0)=r>0.
\end{equation}
On the left exterior interval, monotonicity and $H(-1)<0$ put its zero to
the right of $-1$; on the right exterior interval, $H(0)>0$ puts its zero
to the left of $0$.  All middle zeros lie between poles in $(-1,0)$.
Thus all $n+1$ zeros of $T_rf$ are distinct and belong to $(-1,0)$.

For the remaining cases, first scale $f$ to be monic; this does not change
any zero.  Write its roots as $-s_i$ with $s_i\in[0,1]$.  Choose pairwise
distinct $s_{i,m}\in(0,1)$ converging to $s_i$, and choose
$r_m\in(0,1)$ converging to $r$.  Put
$f_m(x)=\prod_i(x+s_{i,m})$.  By the generic case, all roots of
$T_{r_m}f_m$ lie in the compact interval $[-1,0]$.  These polynomials
converge coefficientwise to $T_rf$, and every one, including the limit, has
degree $n+1$ with leading coefficient $1+n/a>0$.  After normalizing this
leading coefficient, compactness of the ordered root tuples in
$[-1,0]^{n+1}$ and continuity of elementary symmetric functions show that
$T_rf$ also factors into $n+1$ roots in $[-1,0]$.  Finally, if $n=0$ then
$T_rf=(x+r)f$, so the conclusion is immediate.
\end{proof}

\begin{proof}[Proof of Theorem~\ref{thm:main}]
Independent nonzero scalar multiples of $P_1$ and $P_2$ do not affect the
zero sets at issue.  We may therefore write
\[
 \La(P_1)=x+r,\qquad 0\leq r\leq1,
\]
and hence $P_1=e_1+r$.  Put $f=\La(P_2)$.  Lemma~\ref{lem:operator}
identifies $\La(P_1P_2)$ with $T_rf$, and
Lemma~\ref{lem:location} gives the claimed root location.
\end{proof}

For completeness, in the original quadratic case write
$f=(x+s)(x+t)$, where $r,s,t\in[0,1]$.  Formula~\eqref{eq:operator} expands
to
\[
 \La(P_1P_2)=rst+
 \left(r(s+t)+st+\frac{s+t}{a}\right)x+
 \left(r+\frac{(a+1)(s+t)+2}{a}\right)x^2+
 \frac{a+2}{a}x^3.
\]
Theorem~\ref{thm:main} therefore covers the complete semialgebraic parameter
range $a>0$ and $r,s,t\in[0,1]$, including repeated roots and both endpoints.

\section*{Assistance disclosure}

The proof was developed and the manuscript was drafted with AI assistance.
A symbolic-algebra script was used to check the basis-operator identity, the
displayed quadratic-case cubic, and the partial-fraction expansion for sample
degrees.  Those checks were transcription safeguards only; the universal
root-location conclusion follows from the self-contained argument above.
The author takes responsibility for the statements and proof.

\bibliographystyle{plain}
\bibliography{references}

\end{document}
